% ---------------------------------------------------------------------------
% Colour flow and colour chains of our event. Panels (a) and (b) draw the
% hard process and the first initial-state branching in 't Hooft's double-line
% (colour-flow) notation next to the chains the shower keeps; panels (c) and
% (d) show the chains before and after the first two final-state branchings
% (branchings 4 and 5 of the seed-1 shower); panel (e) is colour coherence.
% Every tag, row number and mass is taken from the -v log of
% standalone_parton_showering.py (example/PS/Standalone/standalone_parton_showering.log,
% sections 2 and 4, steps 1 to 5). Input from chapter_02_mother_gives_birth.tex
% inside a figure environment.
% ---------------------------------------------------------------------------
\begin{tikzpicture}[x=1cm,y=1cm,
  every node/.style={font=\footnotesize},
  parton/.style={draw,circle,minimum size=8mm,inner sep=0pt,fill=white,thick,align=center},
  gluon/.style={parton,double,double distance=1.2pt},
  quark/.style={parton},
  beam/.style={draw,rectangle,rounded corners=1.5pt,minimum height=8mm,minimum width=12mm,inner sep=1pt,fill=gray!15,thick,align=center},
  new/.style={fill=yellow!40},
  line/.style={line width=1.8pt},
  tag/.style={font=\scriptsize,fill=white,inner sep=1pt},
  role/.style={font=\scriptsize\itshape,text=black!65,align=center},
  lab/.style={font=\scriptsize,align=center},
  head/.style={font=\small\bfseries},
  % colour-flow lines: one colour index each, arrow along the index
  % (forwards along a quark, backwards along an antiquark)
  cline/.style={line width=1.4pt,rounded corners=2.5pt},
  marks/.style={postaction={decorate},decoration={markings,
     mark=at position #1 with {\arrow[scale=0.9]{Stealth}}}},
  ghost/.style={gray!45,decorate,decoration={coil,aspect=0,segment length=4pt,amplitude=1.4pt}},
  bigarrow/.style={-{Stealth[length=3mm,width=3mm]},gray!60,line width=2.4pt}]
\colorlet{c101}{blue!70!black}
\colorlet{c102}{red!80!black}
\colorlet{c103}{green!55!black}
\colorlet{c104}{orange!90!black}
\colorlet{c105}{violet}
\colorlet{c106}{teal}
\colorlet{c107}{brown!90!black}
% ---------------- (a) the hard process in colour-flow notation --------------------------
\node[head,anchor=west] at (-0.7,2.9) {(a) the hard process $\bar u u\to g^*\to\bbbar$: two colour lines, two dipoles};
\coordinate (V1) at (2.3,0);
\coordinate (V2) at (5.1,0);
\draw[ghost] (V1) -- (V2);
% colour 102: incoming u -> lower edge of the gluon -> b
\draw[cline,c102,marks=0.05,marks=0.5,marks=0.95]
   (0.9,-1.5) -- ($(V1)+(0,-0.13)$) -- ($(V2)+(0,-0.13)$) -- (6.7,-1.5);
% anticolour 101: bbar -> upper edge of the gluon -> incoming ubar (arrow against the fermion)
\draw[cline,c101,marks=0.05,marks=0.5,marks=0.95]
   (6.7,1.5) -- ($(V2)+(0,0.13)$) -- ($(V1)+(0,0.13)$) -- (0.9,1.5);
\node[tag,text=c102] at (1.8,-1.07) {102};
\node[tag,text=c101] at (1.8,1.07) {101};
\node[tag,text=c102] at (5.72,-1.07) {102};
\node[tag,text=c101] at (5.72,1.07) {101};
\node[lab,anchor=east] at (0.8,-1.5) {$u$ \#4\\from the $-z$ proton};
\node[lab,anchor=east] at (0.8,1.5) {$\bar u$ \#3\\from the $+z$ proton};
\node[lab,anchor=west] at (6.8,-1.5) {$b$ \#6, the mother};
\node[lab,anchor=west] at (6.8,1.5) {$\bar b$ \#5};
\node[lab] at (3.7,0.75) {$g^*$: one line of each kind};
\node[lab,text=black!65] at (3.7,-0.75) {$\sqrt{\shat}=396\GeV$};
% the same as chains
\draw[bigarrow] (8.2,0) -- (9.0,0);
\node[role] at (8.6,0.5) {as the shower\\stores it};
\node[beam] (a3) at (10.4,1.0) {$\bar u$ \#3\\[-1pt]{\tiny beam $+z$}};
\node[quark] (a5) at (13.8,1.0) {$\bar b$\\[-2pt]\tiny\#5};
\node[beam] (a4) at (10.4,-1.0) {$u$ \#4\\[-1pt]{\tiny beam $-z$}};
\node[quark] (a6) at (13.8,-1.0) {$b$\\[-2pt]\tiny\#6};
\draw[line,c101] (a3) -- (a5) node[tag,midway,below] {101} node[role,midway,above] {$m_{\mathrm{dip}}=377.9\GeV$};
\draw[line,c102] (a4) -- (a6) node[tag,midway,below] {102} node[role,midway,above] {$m_{\mathrm{dip}}=377.9\GeV$};
\node[role] at (12.5,-1.85) {2 dipoles: the 2 final-state ends radiate (FSR),\\the 2 beam ends evolve backwards (ISR)};
% ---------------- (b) after branching 1 (ISR) -------------------------------------------
\begin{scope}[yshift=-5.4cm]
\node[head,anchor=west] at (-0.7,2.9) {(b) after branching 1 (ISR at $\ptevol=53.9\GeV$): the $\bar u$ came out of a gluon};
\coordinate (V0) at (2.0,1.4);
\coordinate (V1) at (3.4,0);
\coordinate (V2) at (5.4,0);
\draw[ghost] (V1) -- (V2);
\draw[ghost] (0.4,1.4) -- (V0);
% colour 102 as before
\draw[cline,c102,marks=0.05,marks=0.5,marks=0.95]
   (1.7,-1.5) -- ($(V1)+(0,-0.13)$) -- ($(V2)+(0,-0.13)$) -- (7.1,-1.5);
% anticolour 101: bbar -> gluon -> ubar leg -> lower edge of the beam gluon #7
\draw[cline,c101,marks=0.05,marks=0.37,marks=0.66,marks=0.95]
   (7.1,1.5) -- ($(V2)+(0,0.13)$) -- ($(V1)+(0,0.13)$) -- ($(V0)+(0,-0.13)$) -- (0.4,1.27);
% colour 103: upper edge of the beam gluon #7 -> emitted u #8
\draw[cline,c103,marks=0.3,marks=0.8]
   (0.4,1.53) -- ($(V0)+(0,0.13)$) -- (3.5,2.25);
\node[tag,text=c102] at (2.74,-1.07) {102};
\node[tag,text=c102] at (6.07,-1.07) {102};
\node[tag,text=c101] at (6.07,1.07) {101};
\node[tag,text=c101] at (2.45,0.5) {101};
\node[tag,text=c103] at (2.6,2.1) {103};
\node[lab,anchor=east] at (1.6,-1.5) {$u$ \#4 from the $-z$ proton};
\node[lab,anchor=north east] at (1.3,1.15) {$g$ \#7 from the\\$+z$ proton};
\node[lab,anchor=west] at (3.6,2.25) {$u$ \#8 emitted, $45\GeV$};
\node[lab,anchor=west] at (3.05,0.95) {$\bar u$ \#3, spacelike};
\node[lab,anchor=west] at (7.2,-1.5) {$b$ \#10};
\node[lab,anchor=west] at (7.2,1.5) {$\bar b$ \#9};
\node[lab] at (4.4,0.45) {$g^*$};
\draw[bigarrow] (8.5,0) -- (9.3,0);
\node[beam] (b7) at (10.4,1.0) {$g$ \#7\\[-1pt]{\tiny beam $+z$}};
\node[quark] (b8) at (13.8,2.2) {$u$\\[-2pt]\tiny\#8};
\node[quark] (b9) at (13.8,1.0) {$\bar b$\\[-2pt]\tiny\#9};
\node[beam] (b4) at (10.4,-1.0) {$u$ \#4\\[-1pt]{\tiny beam $-z$}};
\node[quark] (b10) at (13.8,-1.0) {$b$\\[-2pt]\tiny\#10};
\draw[line,c103] (b7) -- (b8) node[tag,pos=0.35,below,sloped] {103} node[role,pos=0.68,above,sloped] {$114.6\GeV$};
\draw[line,c101] (b7) -- (b9) node[tag,midway,above] {101} node[role,midway,below] {$420.5\GeV$\\(was $377.9$)};
\draw[line,c102] (b4) -- (b10) node[tag,midway,below] {102} node[role,midway,above] {partner unchanged};
\node[role] at (12.1,-1.85) {3 final-state partons, 3 dipole ends;\\the beam gluon holds two lines};
\end{scope}
% ---------------- notation box, right of (a) and (b) -------------------------------------
\begin{scope}
\fill[gray!8,rounded corners=4pt] (15.6,2.75) rectangle (21.9,-7.3);
\node[head,anchor=west] at (15.8,2.4) {how to read the colour lines};
\draw[cline,c102,marks=0.6] (15.9,1.75) -- (17.1,1.75);
\node[lab,anchor=west] at (17.25,1.75) {quark: one colour, arrow forwards};
\draw[cline,c101,marks=0.6] (17.1,1.15) -- (15.9,1.15);
\node[lab,anchor=west] at (17.25,1.15) {antiquark: one anticolour, arrow backwards};
\draw[ghost] (15.9,0.55) -- (17.1,0.55);
\draw[cline,c103,marks=0.6] (15.9,0.68) -- (17.1,0.68);
\draw[cline,c101,marks=0.6] (17.1,0.42) -- (15.9,0.42);
\node[lab,anchor=west] at (17.25,0.55) {gluon: one of each, two lines};
\node[font=\scriptsize,align=left,anchor=north west,text width=5.6cm] at (15.8,0.05)
  {A quark carries one of $N_c=3$ colours, an antiquark one of three
   anticolours, a gluon one colour and one anticolour. The generator keeps
   only the leading term in $1/N_c$: every diagram becomes a set of lines,
   each carrying one colour index unchanged from where it starts to where
   it ends, and a gluon is a pair of lines with opposite arrows ('t~Hooft's
   double-line notation). The Les Houches record numbers the lines from
   101 upwards; the two colour columns of a parton are the tags of the
   line and the anti-line it holds. A \emph{dipole} is one line with its
   two ends, and the chains on the left list, parton by parton, which
   lines each one holds. Hadronisation will later stretch a string along
   every one of these lines.};
\end{scope}
% ---------------- (c) before branching 4: after the three ISR branchings ------------------
\begin{scope}[yshift=-11.3cm]
\node[head,anchor=west] at (-0.7,3.3) {(c) after the three ISR branchings: 5 final partons, 7 dipole ends};
\node[beam] (b16) at (0,1.0) {$g$ \#16\\[-1pt]{\tiny beam $+z$}};
\node[quark] (u19) at (3.6,2.4) {$u$\\[-2pt]\tiny\#19};
\node[gluon] (g17) at (3.6,1.0) {$g$\\[-2pt]\tiny\#17};
\node[quark] (bb20) at (7.6,1.0) {$\bar b$\\[-2pt]\tiny\#20};
\node[beam] (u11) at (0,-1.2) {$u$ \#11\\[-1pt]{\tiny beam $-z$}};
\node[gluon] (g18) at (3.6,-1.2) {$g$\\[-2pt]\tiny\#18};
\node[quark] (b21) at (7.6,-1.2) {$b$\\[-2pt]\tiny\#21};
\draw[line,c103] (b16) -- (u19) node[tag,midway,above,sloped] {103};
\draw[line,c105] (b16) -- (g17) node[tag,midway,below] {105};
\draw[line,c101] (g17) -- (bb20) node[tag,midway,below] {101} node[role,midway,above] {$m_{\mathrm{dip}}=843\GeV$};
\draw[line,c104] (u11) -- (g18) node[tag,midway,below] {104};
\draw[line,c102] (g18) -- (b21) node[tag,midway,below] {102} node[role,midway,above] {$m_{\mathrm{dip}}=80.6\GeV$};
\node[role,above=1pt of bb20] {radiator of branching 4\\(anticolour 101)};
\node[role,below=1pt of g17] {recoiler of\\branching 4};
\node[role,below=1pt of b21] {radiator of branching 5\\(colour 102)};
\node[role,below=1pt of g18] {recoiler of\\branching 5};
\end{scope}
% ---------------- (d) after branchings 4 and 5 -------------------------------------------
\begin{scope}[yshift=-17.7cm]
\node[head,anchor=west] at (-0.7,3.3) {(d) after branchings 4 and 5: 7 final partons, 11 dipole ends};
\node[beam] (b16) at (0,1.0) {$g$ \#16\\[-1pt]{\tiny beam $+z$}};
\node[quark] (u19) at (3.0,2.4) {$u$\\[-2pt]\tiny\#19};
\node[gluon] (g24) at (3.0,1.0) {$g$\\[-2pt]\tiny\#24};
\node[gluon,new] (g23) at (5.9,1.0) {$g$\\[-2pt]\tiny\#23};
\node[quark] (bb22) at (8.8,1.0) {$\bar b$\\[-2pt]\tiny\#22};
\node[beam] (u11) at (0,-1.2) {$u$ \#11\\[-1pt]{\tiny beam $-z$}};
\node[gluon] (g27) at (3.0,-1.2) {$g$\\[-2pt]\tiny\#27};
\node[gluon,new] (g26) at (5.9,-1.2) {$g$\\[-2pt]\tiny\#26};
\node[quark] (b25) at (8.8,-1.2) {$b$\\[-2pt]\tiny\#25};
\draw[line,c103] (b16) -- (u19) node[tag,midway,above,sloped] {103};
\draw[line,c105] (b16) -- (g24) node[tag,midway,below] {105};
\draw[line,c101] (g24) -- (g23) node[tag,midway,below] {101} node[role,midway,above] {$534\GeV$};
\draw[line,c106] (g23) -- (bb22) node[tag,midway,below] {106} node[role,midway,above] {$69.8\GeV$};
\draw[line,c104] (u11) -- (g27) node[tag,midway,below] {104};
\draw[line,c102] (g27) -- (g26) node[tag,midway,above] {102} node[role,midway,below] {$4.5\GeV$};
\draw[line,c107] (g26) -- (b25) node[tag,midway,above] {107} node[role,midway,below] {$65.7\GeV$};
\node[role,above=1pt of g23] {emitted gluon (status 51), $67\GeV$,\\inserted into line 101};
\node[role,below=1pt of g24] {recoiler's\\new copy (52)};
\node[role,above=1pt of bb22] {$\bar b$'s new\\copy (51)};
\node[role,above=1pt of g26] {emitted gluon (51), $38\GeV$,\\inserted into line 102};
\node[role,below=1pt of g27] {recoiler's\\new copy (52)};
\node[role,below=1pt of b25] {mother's new\\copy (51)};
\end{scope}
% ---------------- (e) colour coherence ------------------------------------------------------
\begin{scope}[xshift=12.9cm,yshift=-13.9cm]
\node[head,anchor=west] at (-1.2,3.6) {(e) what a soft gluon sees};
\coordinate (O) at (0,0);
\coordinate (B) at (20:4.0);
\coordinate (G) at (-32:4.0);
\fill[red!8] (O) -- (46:3.8) arc[start angle=46,end angle=-6,radius=3.8] -- cycle;
\fill[red!8] (O) -- (-6:3.8) arc[start angle=-6,end angle=-58,radius=3.8] -- cycle;
\draw[thick,c107] (O) -- (B) node[pos=1.1,font=\footnotesize] {$b$ \#25};
\draw[thick,c107,decorate,decoration={coil,aspect=0,segment length=4pt,amplitude=1.6pt}] (O) -- (G)
   node[pos=1.12,font=\footnotesize] {$g$ \#26};
\draw[-{Stealth}] (-32:1.5) arc[start angle=-32,end angle=20,radius=1.5];
\node[font=\scriptsize] at (-6:2.0) {$\theta_{bg}$};
\draw[thick,gray!70,decorate,decoration={coil,aspect=0,segment length=3pt,amplitude=1.3pt}] (O) -- (128:3.0)
   node[pos=1.25,font=\scriptsize,align=center,text=black!70] {soft gluon at $\theta\gg\theta_{bg}$:\\sees the pair's total\\charge, $C_F$};
\node[font=\scriptsize,align=left,anchor=north west,text width=5.4cm] at (-1.2,-2.7)
  {Inside the dipole 107 each end radiates mainly within a cone of about $\theta_{bg}$
   (shaded). A gluon at a wider angle cannot resolve the two ends; it sees the colour
   charge of the pair, which is that of the original $b$, $C_F$, not $C_F+C_A$. That is
   colour coherence, and the dipole construction has it built in.};
\end{scope}
\end{tikzpicture}
